%============================================================
% FORMAT: PROGRES PROYEK AKHIR
% Aturan layout khusus Progres PA.
% Di-input SETELAH config/formatting.tex agar bisa override.
%
% Panduan: PASPE-2022
%============================================================

% Geometri: single-sided, margin PASPE-2022
\geometry{
  a4paper,
  top    = 2.54cm,
  bottom = 2.54cm,
  left   = 2.54cm,
  right  = 2.54cm
}

% Spasi Baris (1.15 - lebih kompak dari Buku PA yang pakai 1.5)
\setstretch{1.15}

% Indentasi Paragraf (lebih kecil dari Buku PA)
% Buku PA: 1.27cm. Progres: 0.75cm (~3 karakter Times 12pt)
\setlength{\parindent}{0.75cm}

% Paksa Setiap \chapter ke Halaman Ganjil (opsional)
% Aktif hanya jika \showforceoddpagetrue diset di main_progres.tex.
% Default: off.
\makeatletter
\let\old@chapter\@chapter
\renewcommand{\@chapter}[2][\@empty]{%
  \ifshowforceoddpage\cleartooddpage\fi
  \old@chapter[#1]{#2}%
}
\makeatother

% ============================================================
% OVERRIDE FONT SIZE: 12pt SEMUA (Heading, Caption, TOC, dll)
% ============================================================
\titleformat{\chapter}[display]
{\centering\fontsize{12}{14}\bfseries}
{BAB \thechapter}
{0pt}
{}
\titlespacing*{\chapter}{0pt}{0pt}{12pt}

\titleformat{\section}
{\fontsize{12}{14}\bfseries}
{\thesection}
{1em}
{}
[\SelaraskanIndentAsli{\thesection}]
\titlespacing*{\section}{0pt}{12pt}{12pt}
\titlespacing*{\subsection}{0pt}{10pt}{8pt}

% ── Subsubsection: flush left (mepet margin), tanpa bintang (khusus Progres) ──
% Left = 0pt → subsubsection rata kiri sesuai margin, tidak ikut parindent.
\titleformat{\subsubsection}
{\fontsize{12}{14}\bfseries}
{}
{0pt}
{}
\titlespacing{\subsubsection}{0pt}{3pt}{2pt}

\captionsetup[figure]{font=normalsize}
\captionsetup[table]{font=normalsize}
\SetTblrStyle{caption-tag}{font=\normalsize\itshape}
\SetTblrStyle{caption-text}{font=\normalsize\itshape}

% Override header tabel: abu-abu, bold, center (khusus Progres)
% Tidak diset di config/formatting.tex agar Buku PA tidak ikut.
\SetTblrInner{row{1} = {bg=gray!20, font=\bfseries, halign=c}}
\SetTblrInner[longtblr]{row{1} = {bg=gray!20, font=\bfseries, halign=c}}

\renewcommand{\cfttoctitlefont}{\hfill\fontsize{12}{14}\bfseries\MakeUppercase}
\renewcommand{\cftloftitlefont}{\hfill\fontsize{12}{14}\bfseries\MakeUppercase}
\renewcommand{\cftlottitlefont}{\hfill\fontsize{12}{14}\bfseries\MakeUppercase}

\renewenvironment{kutipanpanjang}{%
  \begin{list}{}{%
      \setlength{\leftmargin}{10mm}
      \setlength{\rightmargin}{10mm}
      \setlength{\topsep}{6pt}
      \setlength{\itemsep}{0pt}
      \setlength{\parsep}{0pt}
    }
  \item\fontsize{12}{14}\selectfont\ignorespaces
  }{%
    \unskip
  \end{list}
}
